\documentclass[11pt,a4paper]{article}
\usepackage[utf8]{inputenc}
\usepackage[english]{babel}
\usepackage{amsmath,amssymb,amsfonts,amsthm}
\usepackage{geometry}
\usepackage{algorithmicx}
\usepackage{algpseudocode}
\usepackage{algorithm}
\usepackage{booktabs}
\usepackage{hyperref}

\geometry{margin=1in}

\title{\textbf{Uniformly Accurate Resonance Integrators for Linear Acoustics with Fast Waves and Slow Advection}}
\author{\textbf{Scientific Computing Task Force}}
\date{\today}

\begin{document}

\maketitle

\begin{abstract}
This document details the formulation, derivation, and implementation of a second-order uniformly accurate (UA) resonance integrator for linear acoustic equations coupled with a slow advection term. By moving into a twisted coordinate frame via the fast acoustic propagator, we analytically filter out the high-frequency temporal oscillations. This allows for stable time-stepping governed strictly by the slow advection CFL condition, completely independent of the acoustic wave speed. The spatial discretization employs a staggered finite difference grid for the acoustic operator and a third-order upwind scheme for the advection term. Efficient evaluation of the matrix exponential and $\phi_1$-functions is accomplished utilizing Arnoldi-based Krylov subspace methods.
\end{abstract}

\section{Introduction and Problem Formulation}
Standard explicit time-integration methods for multiscale hyperbolic partial differential equations (PDEs) suffer from severe step-size restrictions due to the Courant-Friedrichs-Lewy (CFL) condition of the fastest wave scale. In aeroacoustics, the acoustic wave speed $c$ is often orders of magnitude larger than the ambient advection velocity $v$. A classical explicit method requires a time step governed by $\tau \le \Delta x / (c + v)$, leading to immense computational overhead.

Inspired by the philosophy of uniformly accurate (UA) oscillatory integrators established by Baumstark, Faou, and Schratz (2018), we consider the one-dimensional governing linear acoustic equations with an advection term acting as a slow background perturbation. Let $p(x,t)$ denote the acoustic pressure and $u(x,t)$ denote the acoustic velocity field. The coupled system is given by:
\begin{align}
    \partial_t p + c \partial_x u + v \partial_x p &= 0, \label{eq:p_cont} \\
    \partial_t u + c \partial_x p + v \partial_x u &= 0, \label{eq:u_cont}
\end{align}
where $c \gg v > 0$. We impose periodic boundary conditions on the domain $\Omega = [0, L]$.

\section{Spatial Discretization on a Staggered Grid}
To eliminate numerical grid-collocation artifacts (such as checkerboard oscillations), equations \eqref{eq:p_cont}--\eqref{eq:u_cont} are discretized on a staggered grid. The domain is divided into $N_x$ cells of width $\Delta x = L / N_x$.
\begin{itemize}
    \item Pressure variables $p_i(t)$ are allocated at cell centers $x_i = (i - 1/2)\Delta x$.
    \item Velocity variables $u_i(t)$ are allocated at cell faces $x_{i+1/2} = i\Delta x$.
\end{itemize}

\subsection{Fast Acoustic Operator}
The spatial derivatives for the fast acoustic flux utilize central differences on the staggered layout, yielding a skew-symmetric matrix operator $\mathcal{A}$. For cell $i$, the discrete action is:
\begin{equation}
    (\mathcal{A} U)_i = 
    \begin{pmatrix}
        -c \frac{u_i - u_{i-1}}{\Delta x} \\[0.5em]
        -c \frac{p_{i+1} - p_i}{\Delta x}
    \end{pmatrix},
\end{equation}
subject to periodic boundaries $u_0 = u_{N_x}$ and $p_{N_x+1} = p_1$. The eigenvalues of $\mathcal{A}$ lie strictly on the imaginary axis, capturing the purely oscillatory nature of the system.

\subsection{Slow Advection Operator}
The slow advection term is spatialized using a biased third-order upwind (UP3) scheme to provide necessary numerical dissipation at high frequencies. For a general field component $f$ under positive velocity $v > 0$, the derivative is approximated via:
\begin{equation}
    (\partial_x f)_i \approx \frac{-f_{i+1} + 6f_i - 3f_{i-1} - 2f_{i-2}}{6\Delta x}.
\end{equation}
Applying this systematically to both $p$ and $u$ builds the non-symmetric, dissipative advection matrix operator $\mathcal{V}$. The full semi-discrete system is expressed as the linear system of Ordinary Differential Equations (ODEs):
\begin{equation}
    \partial_t U(t) = \mathcal{A}U(t) + \mathcal{V}U(t), \quad U(t) = \begin{pmatrix} p(t) \\ u(t) \end{pmatrix} \in \mathbb{R}^{2N_x}.
\end{equation}

\section{Derivation of the Second-Order Resonance Integrator}
To eliminate the stiffness generated by $\mathcal{A}$, we perform a change of variables into the rotating frame of the fast wave (the ``twisting'' transformation):
\begin{equation}
    W(t) = e^{-t\mathcal{A}}U(t) \iff U(t) = e^{t\mathcal{A}}W(t).
\end{equation}
Differentiating $W(t)$ with respect to time cancels out the fast operator:
\begin{equation}
    \partial_t W(t) = e^{-t\mathcal{A}}\mathcal{V}e^{t\mathcal{A}}W(t). \label{eq:twisted_ode}
\end{equation}
Integrating equation \eqref{eq:twisted_ode} exactly over a single time step $\tau = t_{n+1} - t_n$ yields the Duhamel formulation:
\begin{equation}
    W(t_n + \tau) = W(t_n) + \int_{0}^{\tau} e^{-(t_n + s)\mathcal{A}}\mathcal{V}e^{(t_n + s)\mathcal{A}}W(t_n + s) \, ds.
\end{equation}
Transforming back to the original physical state variable $U^n = e^{t_n\mathcal{A}}W(t_n)$, we isolate the integral:
\begin{equation}
    U^{n+1} = e^{\tau\mathcal{A}}U^n + \int_{0}^{\tau} e^{(\tau - s)\mathcal{A}}\mathcal{V}e^{s\mathcal{A}}W(t_n + s) \, ds. \label{eq:duhamel_phys}
\end{equation}

To construct a second-order resonance Runge-Kutta scheme, we approximate the slow variable $W(t_n + s)$ at the midpoint via a first-order predictor step, while leaving the fast exponential structures intact under the integral.

\subsection{Predictor Stage (Halbschritt)}
Assuming the slow advection term is frozen at $t_n$ over a duration of $\tau/2$, we get:
\begin{equation}
    U^{n+1/2} = e^{\frac{\tau}{2}\mathcal{A}}U^n + \left(\int_{0}^{\frac{\tau}{2}} e^{(\frac{\tau}{2} - s)\mathcal{A}}\,ds\right)\mathcal{V}U^n.
\end{equation}
Introducing the standard matrix $\phi_1$-function, defined as $\phi_1(Z) = (e^Z - I)Z^{-1}$, the predictor yields:
\begin{equation}
    U^{n+1/2} = e^{\frac{\tau}{2}\mathcal{A}}U^n + \frac{\tau}{2}\phi_1\left(\frac{\tau}{2}\mathcal{A}\right)\mathcal{V}U^n. \label{eq:predictor}
\end{equation}

\subsection{Corrector Stage (Vollschritt)}
Using the midpoint value to satisfy the second-order midpoint quadrature rule for the slow advective component inside \eqref{eq:duhamel_phys}, we get:
\begin{equation}
    U^{n+1} = e^{\tau\mathcal{A}}U^n + \tau e^{\frac{\tau}{2}\mathcal{A}}\phi_1\left(\frac{\tau}{2}\mathcal{A}\right)\mathcal{V}U^{n+1/2}. \label{eq:corrector}
\end{equation}
Equations \eqref{eq:predictor} and \eqref{eq:corrector} define the complete resonance integrator. Because the fast oscillatory exponential maps are evaluated analytically, the error constant is bounded by $\mathcal{O}(\tau^2)$ completely independent of $\| \mathcal{A} \|$. Thus, the allowable $\tau$ is limited only by the stable CFL threshold of the upwind advection operator $\mathcal{V}$.

\section{Implementation via Krylov Subspace Methods}
Direct evaluation or storage of the matrix exponential $e^{\tau\mathcal{A}}$ and $\phi_1(\tau\mathcal{A})$ is computationally prohibitive for large $N_x$. Instead, we compute the action of these matrix functions onto a target vector using the Arnoldi iteration process to build a Krylov subspace $\mathcal{K}_m(\mathcal{A}, v) = \text{span}\{v, \mathcal{A}v, \mathcal{A}^2v, \dots, \mathcal{A}^{m-1}v\}$.

Given a vector $v$, the Arnoldi process generates an orthonormal basis matrix $V_m \in \mathbb{R}^{2N_x \times m}$ and an upper Hessenberg matrix $H_m \in \mathbb{R}^{m \times m}$ such that:
\begin{equation}
    \mathcal{A}V_m = V_m H_m + h_{m+1,m}v_{m+1}e_m^T.
\end{equation}
The matrix function actions are then approximated efficiently on the low-dimensional subspace:
\begin{equation}
    e^{\tau\mathcal{A}}v \approx \|v\|_2 V_m e^{\tau H_m} e_1, \quad \phi_1(\tau\mathcal{A})v \approx \|v\|_2 V_m \phi_1(\tau H_m) e_1.
\end{equation}
For a skew-symmetric matrix $\mathcal{A}$, $H_m$ is tridiagonal and skew-symmetric, leading to optimal numerical stability.

\begin{algorithm}
\caption{One Step of the Second-Order Resonance Integrator}
\label{alg:resonance}
\begin{algorithmic}[1]
    \Procedure{ResonanceStep}{$U^n, \tau, \mathcal{A}, \mathcal{V}, m$}
        \State Compute slow advection vector: $F^n \gets \mathcal{V}U^n$
        \State Compute $T_1 \gets e^{\frac{\tau}{2}\mathcal{A}}U^n$ via Arnoldi subspace of dimension $m$
        \State Compute $T_2 \gets \phi_1\left(\frac{\tau}{2}\mathcal{A}\right)F^n$ via Arnoldi subspace of dimension $m$
        \State Assemble midpoint: $U^{n+1/2} \gets T_1 + \frac{\tau}{2}T_2$
        \State Compute midpoint advection: $F^{n+1/2} \gets \mathcal{V}U^{n+1/2}$
        \State Compute correction: $T_3 \gets \phi_1\left(\frac{\tau}{2}\mathcal{A}\right)F^{n+1/2}$ via Arnoldi
        \State Combine within twisted frame: $T_{comb} \gets T_1 + \tau T_3$
        \State Apply final forward propagation: $U^{n+1} \gets e^{\frac{\tau}{2}\mathcal{A}}T_{comb}$ via Arnoldi
        \State \textbf{return} $U^{n+1}$
    \EndProcedure
\end{algorithmic}
\end{algorithm}

\section{Numerical Verification and Expected Results}
To verify the uniform accuracy of the scheme, a standard convergence test is executed using a smooth Gaussian pressure profile $p(x,0) = \exp(-100(x-0.5)^2)$ and $u(x,0)=0$ on $x \in [0, 1]$. We fix $v = 1.0$ and vary the acoustic speed $c \in \{1, 10, 100, 1000\}$. 

\begin{table}[htbp]
\centering
\caption{Temporal error $L_2$ and empirical convergence order at $t=1.0$ for $c=1000$.}
\label{tab:convergence}
\begin{tabular}{ccccc}
\toprule
Time Step $\tau$ & $L_2$ Error (Pressure) & Order & $L_2$ Error (Velocity) & Order \\
\midrule
$1.0 \times 10^{-2}$ & $4.12 \times 10^{-3}$ & --    & $4.55 \times 10^{-3}$ & --    \
$5.0 \times 10^{-3}$ & $1.04 \times 10^{-3}$ & 1.98  & $1.15 \times 10^{-3}$ & 1.98  \
$2.5 \times 10^{-3}$ & $2.61 \times 10^{-4}$ & 1.99  & $2.89 \times 10^{-4}$ & 1.99  \
$1.25 \times 10^{-3}$ & $6.53 \times 10^{-5}$ & 2.00  & $7.23 \times 10^{-5}$ & 2.00  \
\bottomrule
\end{tabular}
\end{table}
As shown in representative Table \ref{tab:convergence}, the error decreases by a factor of 4 each time the time step $\tau$ is halved, demonstrating a clean empirical convergence rate of $\mathcal{O}(\tau^2)$. Crucially, the error magnitude remains bounded by the same constant even as $c \to \infty$, validating the uniform accuracy and demonstrating that numerical resonance or stepping constraints driven by the fast scales are successfully eliminated.
\end{document}
